%!TEX program=xelatex
\documentclass[11pt,a4paper]{report}

% =================================================
% ЯЗЫК И ШРИФТЫ
% =================================================
\usepackage{fontspec}
\usepackage{polyglossia}
\setdefaultlanguage{russian}
\setotherlanguage{english}
\IfFontExistsTF{Times New Roman}{\setmainfont{Times New Roman}}{\setmainfont{Liberation Serif}}
\IfFontExistsTF{Arial}{\setsansfont{Arial}}{\setsansfont{Liberation Sans}}
\IfFontExistsTF{Menlo}{\setmonofont{Menlo}}{\setmonofont{DejaVu Sans Mono}}

% =================================================
% ГЕОМЕТРИЯ СТРАНИЦЫ
% =================================================
\usepackage[
    a4paper,
    left=1.7cm,
    right=1.7cm,
    top=1.8cm,
    bottom=1.8cm,
    headsep=0.7cm,
    footskip=1cm
]{geometry}

% =================================================
% МАТЕМАТИКА
% =================================================
\usepackage{amsmath,amssymb,amsfonts,mathtools}
\usepackage{mathrsfs}
\usepackage{bm}
\usepackage{gensymb}

% =================================================
% ГРАФИКА И ТАБЛИЦЫ
% =================================================
\usepackage{graphicx}
\usepackage{caption}
\usepackage{subcaption}
\usepackage{float}
\usepackage{booktabs}
\usepackage{multirow}
\usepackage{array}
\usepackage{longtable}
\usepackage{tabularx}
\usepackage{comment}

% =================================================
% TIKZ / PGFPLOTS
% =================================================
\usepackage{tikz}
\usetikzlibrary{patterns,arrows.meta,calc}
\usepackage{pgfplots}
\pgfplotsset{compat=1.18}

% =================================================
% ТИПОГРАФИКА
% =================================================
\usepackage{microtype}
\usepackage{ragged2e}
\usepackage{indentfirst}
\usepackage{setspace}
\setstretch{1.08}
\setlength{\parindent}{0.7cm}
\setlength{\parskip}{0pt}
\setlength{\abovedisplayskip}{8pt}
\setlength{\belowdisplayskip}{8pt}
\setlength{\abovedisplayshortskip}{5pt}
\setlength{\belowdisplayshortskip}{5pt}

% =================================================
% СПИСКИ
% =================================================
\usepackage{enumitem}
\setlist[itemize]{topsep=4pt,itemsep=1pt,parsep=0pt,leftmargin=1.8em}
\setlist[enumerate]{topsep=4pt,itemsep=1pt,parsep=0pt,leftmargin=2em}

% =================================================
% ГИПЕРССЫЛКИ
% =================================================
\usepackage[
    unicode=true,
    colorlinks=true,
    linkcolor=blue,
    citecolor=blue,
    urlcolor=blue,
    linktocpage=true
]{hyperref}
\hypersetup{
    pdftitle={Конспект по математической физике},
    pdfauthor={Яков Овсянников, Александр Лозовой}
}
\usepackage{aliascnt}
\usepackage[nameinlink,noabbrev]{cleveref}

% =================================================
% СТРАНИЦЫ
% =================================================
% Номера страниц снизу; титульная страница отдельно получает \thispagestyle{empty}.
\pagestyle{plain}

% =================================================
% ОФОРМЛЕНИЕ ГЛАВ И РАЗДЕЛОВ
% =================================================
\usepackage{titlesec}
\titleformat{\chapter}[hang]
    {\normalfont\LARGE\bfseries}
    {Глава \thechapter.}
    {0.6em}
    {}
\titlespacing*{\chapter}{0pt}{-20pt}{18pt}

\titleformat{\section}
    {\normalfont\Large\bfseries}
    {\thesection.}
    {0.5em}
    {}
\titlespacing*{\section}{0pt}{2.0ex}{1.0ex}

\titleformat{\subsection}
    {\normalfont\large\bfseries}
    {\thesubsection.}
    {0.5em}
    {}
\titlespacing*{\subsection}{0pt}{1.7ex}{0.7ex}

\titleformat{\subsubsection}
    {\normalfont\normalsize\bfseries}
    {\thesubsubsection.}
    {0.5em}
    {}

% =================================================
% ОГЛАВЛЕНИЕ
% =================================================
\usepackage{tocloft}
\renewcommand{\cfttoctitlefont}{\hfill\Large\bfseries}
\renewcommand{\cftaftertoctitle}{\hfill}
\renewcommand{\cftchapfont}{\bfseries}
\renewcommand{\cftchappagefont}{\bfseries}
\setlength{\cftbeforechapskip}{0.35em}

% =================================================
% НУМЕРАЦИЯ СТРУКТУРЫ
% =================================================
% Глава: 1; раздел: 1.1; подраздел: 1.1.1.
\renewcommand{\thesection}{\thechapter.\arabic{section}}
\renewcommand{\thesubsection}{\thesection.\arabic{subsection}}
\renewcommand{\thesubsubsection}{\thesubsection.\arabic{subsubsection}}

% Формулы нумеруются внутри разделов: (1.2.1), (1.2.2), ...
\numberwithin{equation}{section}
% Рисунки и таблицы — внутри главы: 1.1, 1.2, ...
\numberwithin{figure}{chapter}
\numberwithin{table}{chapter}
\setcounter{secnumdepth}{3}
\setcounter{tocdepth}{3}

% =================================================
% ПОДПИСИ К РИСУНКАМ И ТАБЛИЦАМ
% =================================================
\captionsetup{labelfont=bf,font=small,labelsep=period}

% =================================================
% ЕДИНАЯ НУМЕРАЦИЯ МАТЕМАТИЧЕСКИХ БЛОКОВ
% =================================================
% Все определения, теоремы, леммы, утверждения, примеры,
% упражнения и замечания идут одной последовательностью внутри section.
\newcounter{statementctr}[section]
\renewcommand{\thestatementctr}{\thesection.\arabic{statementctr}}

\newaliascnt{theoremctr}{statementctr}
\newaliascnt{lemmactr}{statementctr}
\newaliascnt{propctr}{statementctr}
\newaliascnt{defctr}{statementctr}
\newaliascnt{remarkctr}{statementctr}
\newaliascnt{examplectr}{statementctr}
\newaliascnt{exectr}{statementctr}

\aliascntresetthe{theoremctr}
\aliascntresetthe{lemmactr}
\aliascntresetthe{propctr}
\aliascntresetthe{defctr}
\aliascntresetthe{remarkctr}
\aliascntresetthe{examplectr}
\aliascntresetthe{exectr}

\crefname{theoremctr}{теорема}{теоремы}
\Crefname{theoremctr}{Теорема}{Теоремы}
\crefname{lemmactr}{лемма}{леммы}
\Crefname{lemmactr}{Лемма}{Леммы}
\crefname{propctr}{утверждение}{утверждения}
\Crefname{propctr}{Утверждение}{Утверждения}
\crefname{remarkctr}{замечание}{замечания}
\Crefname{remarkctr}{Замечание}{Замечания}
\crefname{defctr}{определение}{определения}
\Crefname{defctr}{Определение}{Определения}
\crefname{exectr}{упражнение}{упражнения}
\Crefname{exectr}{Упражнение}{Упражнения}
\crefname{examplectr}{пример}{примеры}
\Crefname{examplectr}{Пример}{Примеры}

% Флаг нужен только для совместимости со старыми \lbpa/\lbpn.
\newif\ifinmathphysstatement
\inmathphysstatementfalse

% Старые \lbpa оставлены, чтобы старый текст продолжал компилироваться.
% Теперь они используют общий счётчик математических утверждений.
\newcommand{\lbpa}[2]{%
    \refstepcounter{statementctr}%
    \label{#1}%
    \noindent #2 \hfill(\thestatementctr)%
    \par\medskip
}
\newcommand{\lbpn}[2]{%
    \ifinmathphysstatement
        \label{#1}%
        \noindent #2%
    \else
        \refstepcounter{equation}%
        \label{#1}%
        \noindent #2 \hfill(\theequation)%
    \fi
    \par\medskip
}

% =================================================
% ЦВЕТНЫЕ БЛОКИ
% =================================================
\usepackage[most]{tcolorbox}

% Спокойная палитра, близкая по логике к main.tex.
\definecolor{defcolor}{RGB}{46,125,95}
\definecolor{theoremcolor}{RGB}{55,95,160}
\definecolor{lemmacolor}{RGB}{105,82,150}
\definecolor{propcolor}{RGB}{55,120,135}
\definecolor{examplecolor}{RGB}{180,90,75}
\definecolor{exercisecolor}{RGB}{40,125,155}
\definecolor{remarkcolor}{RGB}{175,125,45}
\definecolor{proofcolor}{RGB}{105,105,105}

\tcbset{
    mathphysblock/.style={
        enhanced,
        breakable,
        boxrule=0.65pt,
        arc=1.5mm,
        left=7pt,
        right=7pt,
        top=6pt,
        bottom=6pt,
        before skip=9pt,
        after skip=9pt
    }
}

\newenvironment{mydef}[1][]{%
    \refstepcounter{defctr}%
    \inmathphysstatementtrue
    \begin{tcolorbox}[mathphysblock,colframe=defcolor,colback=defcolor!3,#1]
    \textbf{\color{defcolor}Определение \thedefctr.}\par\smallskip
}{%
    \end{tcolorbox}
    \inmathphysstatementfalse
}

\newenvironment{mytheorem}[1][]{%
    \refstepcounter{theoremctr}%
    \inmathphysstatementtrue
    \begin{tcolorbox}[mathphysblock,colframe=theoremcolor,colback=theoremcolor!3,boxrule=0.8pt,#1]
    \textbf{\color{theoremcolor}Теорема \thetheoremctr.}\par\smallskip
}{%
    \end{tcolorbox}
    \inmathphysstatementfalse
}

\newenvironment{mylemma}[1][]{%
    \refstepcounter{lemmactr}%
    \inmathphysstatementtrue
    \begin{tcolorbox}[mathphysblock,colframe=lemmacolor,colback=lemmacolor!3,#1]
    \textbf{\color{lemmacolor}Лемма \thelemmactr.}\par\smallskip
}{%
    \end{tcolorbox}
    \inmathphysstatementfalse
}

\newenvironment{myprop}[1][]{%
    \refstepcounter{propctr}%
    \inmathphysstatementtrue
    \begin{tcolorbox}[mathphysblock,colframe=propcolor,colback=propcolor!3,#1]
    \textbf{\color{propcolor}Утверждение \thepropctr.}\par\smallskip
}{%
    \end{tcolorbox}
    \inmathphysstatementfalse
}

\newenvironment{myexample}[1][]{%
    \refstepcounter{examplectr}%
    \inmathphysstatementtrue
    \begin{tcolorbox}[mathphysblock,colframe=examplecolor,colback=examplecolor!3,#1]
    \textbf{\color{examplecolor}Пример \theexamplectr.}\par\smallskip
}{%
    \end{tcolorbox}
    \inmathphysstatementfalse
}

\newenvironment{myexe}[1][]{%
    \refstepcounter{exectr}%
    \inmathphysstatementtrue
    \begin{tcolorbox}[mathphysblock,colframe=exercisecolor,colback=exercisecolor!3,#1]
    \textbf{\color{exercisecolor}Упражнение \theexectr.}\par\smallskip
}{%
    \end{tcolorbox}
    \inmathphysstatementfalse
}

\newenvironment{myremark}[1][]{%
    \refstepcounter{remarkctr}%
    \inmathphysstatementtrue
    \begin{tcolorbox}[mathphysblock,colframe=remarkcolor,colback=remarkcolor!4,#1]
    \textbf{\color{remarkcolor}Замечание \theremarkctr.}\par\smallskip
}{%
    \end{tcolorbox}
    \inmathphysstatementfalse
}

\newenvironment{myproof}{%
    \begin{tcolorbox}[mathphysblock,colframe=proofcolor!65,colback=proofcolor!2,boxrule=0.5pt]
    \textit{\textbf{Доказательство.}}\par\smallskip
}{%
    \hfill$\square$
    \end{tcolorbox}
}

% -------------------------------------------------
% ПОЛЕЗНЫЕ КОМАНДЫ
% -------------------------------------------------
\newcommand{\R}{\mathbb{R}}
\newcommand{\N}{\mathbb{N}}
\newcommand{\Z}{\mathbb{Z}}
\newcommand{\Q}{\mathbb{Q}}
\newcommand{\C}{\mathbb{C}}
\newcommand{\E}{\mathbb{E}}
\newcommand{\D}{\mathcal{D}}
\newcommand{\B}{\mathcal{B}}
\newcommand{\F}{\mathcal{F}}
\newcommand{\LL}{\mathscr{L}}
\renewcommand{\L}{\mathcal{L}}
\newcommand{\circled}[1]{\tikz[baseline=(char.base)] \node[draw,circle,inner sep=1pt,minimum size=1.4em] (char) {#1};}

\newcommand{\sign}{\operatorname{sign}}
\renewcommand{\span}{\operatorname{span}}
\newcommand{\cl}{\operatorname{cl}}
\newcommand{\spec}{\operatorname{spec}}
\newcommand{\dom}{\operatorname{Dom}}
\newcommand{\Alt}{\operatorname{Alt}}
\newcommand{\supp}{\operatorname{supp}}
\newcommand{\arcsh}{\operatorname{arcsh}}
\newcommand{\arcch}{\operatorname{arcch}}
\newcommand{\arctanh}{\operatorname{arctanh}}
\newcommand{\arccth}{\operatorname{arccth}}

\newcommand{\DR}{\mathcal{D}(\R)}
\newcommand{\SR}{\mathcal{S}(\R)}
\newcommand{\SSR}{\mathcal{S}'(\R)}
\newcommand{\dR}{\mathcal{D}(\R^n)}
\newcommand{\ddR}{\mathcal{D}'(\R^n)}
\newcommand{\DDR}{\mathcal{D}'(\R)}
\newcommand{\Cinf}{\text{C}^{\infty}}

\newcommand{\p}{\partial}
\newcommand{\dd}{\mathop{}\!\mathrm{d}}
\newcommand{\ee}{\mathrm{e}}
\newcommand{\ii}{\mathrm{i}}
\newcommand{\grad}{\Vec{\nabla}}

\newcommand{\abs}[1]{\left|#1\right|}
\newcommand{\norm}[1]{\left\|#1\right\|}
\newcommand{\avg}[1]{\left\langle #1 \right\rangle}
\newcommand{\inner}[1]{\left\langle #1 \right\rangle}
\newcommand{\podst}[1]{\left[ #1 \right]}
\newcommand{\skobka}[1]{\left( #1 \right)}
\newcommand{\paren}[1]{\left(#1\right)}
\newcommand{\brackets}[1]{\left[#1\right]}
\newcommand{\set}[1]{\left\{#1\right\}}

\newcommand{\Pf}[1]{\mathrm{P}\dfrac{1}{#1}}
\newcommand{\PP}{\mathrm{P}\dfrac{1}{x}}
\newcommand{\px}[1]{\dfrac{1}{#1 \pm \mathrm{i}0}}
\newcommand{\Y}{\mathscr{Y}}

\newcommand{\INT}{\int\limits_{-\infty}^{+\infty}}
\newcommand{\intpi}{\int\limits_{-\pi}^{+\pi}}
\newcommand{\sumn}{\sum_{n=0}^{\infty}}
\newcommand{\summ}{\sum_{m=0}^{\infty}}
\newcommand{\sumk}{\sum_{k=0}^{\infty}}
\newcommand{\SUMN}{\sum\limits_{n=1}^{\infty}}
\newcommand{\SUMM}{\sum\limits_{m=1}^{\infty}}
\newcommand{\SUMK}{\sum\limits_{k=1}^{\infty}}
\newcommand{\sumin}{\sum_{i=1}^{n}}
\newcommand{\sumkn}{\sum_{k=1}^{n}}
\newcommand{\SuM}{\sum\limits_{k=1}^{m}}
\newcommand{\disp}{\displaystyle}

\newcommand{\smallo}[2]{\underset{#1}{\overline{o}\left(#2\right)}}
\newcommand{\bigo}[2]{\mathop{O}\limits_{#1}\!\left(#2\right)}
\newcommand{\bigO}{\underline{O}}

\newcommand{\vect}[1]{\mathbf{#1}}
\newcommand{\Let}{\sqsupset}
\newcommand{\ocup}{\overset{\circ}{\cup}}
\newcommand{\dopis}{\begin{center}\textbf{Дописать}\end{center}}
\newcommand{\ubf}[1]{\noindent\textbf{\underline{#1}}}
\newcommand{\theme}[1]{\noindent\textbf{\underline{#1}}\vskip2mm}
\newcommand{\RomanNumeralCaps}[1]{\MakeUppercase{\romannumeral #1}}
\newcounter{strcounter}
\newcommand{\numberedline}[1]{%
    \refstepcounter{strcounter}%
    \makebox[0pt][r]{\thestrcounter\hspace{1em}}%
    \label{#1}%
}

\renewcommand{\phi}{\varphi}
\renewcommand{\epsilon}{\varepsilon}
\renewcommand{\kappa}{\varkappa}

\begin{document}


% =================================================
% ОБЛОЖКА
% =================================================
\begin{titlepage}
\thispagestyle{empty}
\centering

\vspace*{0.8cm}

{\Huge\bfseries Конспект по математической физике \par}
\vspace{0.4cm}

{\Large Осенний семестр 2026\par}

\vspace{0.6cm}

%\includegraphics[width=0.6\textwidth]{pics/11.png}

\vspace{0.4cm}

\vfill

{\large Лекторы\par}
\vspace{0.2cm}
{\Large Филонов Николай Дмитриевич\par}



{\large \href{https://www.youtube.com/playlist?list=PLY00npmGiNp8}
{Плейлист лекций}}
\vspace{0.2cm}
{\normalsize

\par}

\vspace{0.7cm}

{\large Над файлом работали\par}
\vspace{0.2cm}
{\normalsize
Яков Овсянников(ТГ:@Yasha2707)\par }
{\normalsize
Александр Лозовой(ТГ:@sdlhr)\par}

\vspace{1.2cm}

{\large Университет ИТМО\par}

\end{titlepage}
% =================================================
% РЕКОМЕНДУЕМАЯ ЛИТЕРАТУРА
% =================================================

\clearpage
\section*{Рекомендуемая литература}

\begin{enumerate}
    \item Буслаев В.\,С.
    \textit{Вариационное исчисление}.

    \item Гельфанд И.\,М., Фомин С.\,В.
    \textit{Вариационное исчисление}.

    \item Эльсгольц Л.\,Э.
    \textit{Вариационное исчисление}.

    \item Смирнов В.\,И.
    \textit{Курс высшей математики}. Том 4.
\end{enumerate}

\clearpage

% =================================================
% ОГЛАВЛЕНИЕ
% =================================================
\tableofcontents
\clearpage



\section*{\S 0. Разбиение единицы}
\addcontentsline{toc}{section}{\S 0. Разбиение единицы}
С прошлого семетсра у нас остался долг по обсуждению сюжета про разбиение единицы. В этом разделе мы разберемся с этим вопросом.\\
Рассмотрим функцию
\[
\omega(x)=
\begin{cases}
C\exp\!\left(-\dfrac{1}{1-|x|^2}\right), & |x|<1,\\[2mm]
0, & |x|\ge 1,
\end{cases}
\]
где нормировочная константа $C$ выбрана так, что
\[
\omega(x)>0,
\qquad
\int_{\mathbb R^n}\omega(x)\,\dd x=1.
\]
Тогда
\[
\omega\in C_0^\infty(\mathbb R^n),
\qquad
\operatorname{supp}\omega=\overline{B_1(0)}.
\]

Для любого $\rho>0$ положим
\[
\omega_\rho(x)=\frac{1}{\rho^n}\omega\!\left(\frac{x}{\rho}\right).
\]
Тогда
\[
\omega_\rho\in C_0^\infty(\mathbb R^n),
\qquad
\omega_\rho>0,
\qquad
\operatorname{supp}\omega_\rho=\overline{B_\rho(0)},
\qquad
\int_{\mathbb R^n}\omega_\rho(x)\,\dd x=1.
\]

\begin{mylemma}
Пусть $\Omega$ и $U$ --- открытые множества и
$
\overline\Omega\subset U.
$
Тогда существует функция $\eta$, 
\[
\eta\in C_0^\infty(\mathbb R^n),
\qquad
0\le \eta(x)\le 1\quad \forall x,
\]
такая, что
\[
\eta\big|_{\overline\Omega}=1,
\qquad
\operatorname{supp}\eta\subset U.
\]
\end{mylemma}

\begin{myproof}
Положим
\[
\rho:=\frac14\operatorname{dist}(\overline\Omega,\partial U)>0.
\]
Введём множество
\[
A:=\left\{x\in\mathbb R^n:\operatorname{dist}(x,\overline\Omega)<2\rho\right\}.
\]
Тогда $A$ открыто и
$
\overline\Omega\subset A,
\qquad
\overline A\subset U.
$

Положим
\[
\eta(x):=\int_A\omega_\rho(x-y)\,\dd y
=(\omega_\rho*\chi_A)(x).
\]
Имеем
\[
\eta\in C^\infty,
\qquad
0\le \eta(x)\le 1.
\]


Если $x\in\overline\Omega$, то
\[
B_\rho(x)\subset A \Rightarrow
\eta(x)
=
\int_{B_\rho(x)}\omega_\rho(x-y)\,\dd y
=1.
\]

Если $x\notin U$, то
\[
B_\rho(x)\cap A=\varnothing \Rightarrow
\eta(x)=0.
\]
\end{myproof}

\begin{mylemma}
Пусть $\Omega$ --- открытая ограниченная область,
\[
\overline\Omega\subset\bigcup_{j=1}^{N}O_j,
\qquad O_j\text{ --- открытые}.
\]
Тогда существуют открытые множества
\[
\Omega_j,\qquad j=0,1,\ldots,N,
\]
такие, что
\[
\overline\Omega\subset\bigcup_{j=0}^{N}\Omega_j,
\qquad
\overline{\Omega_j}\subset O_j,
\quad j=1,\ldots,N,
\]
и
\[
\overline{\Omega_0}\cap\overline\Omega=\varnothing.
\]
\end{mylemma}

\begin{myproof}
Выберем куб $Q$ так, чтобы рассматриваемое покрытие содержалось внутри некоторого большего открытого куба $\widetilde Q$.
Последовательно строим открытые множества $\Omega_j$ так, чтобы
\[
\overline{\Omega_j}\subset O_j,
\qquad j=1,\ldots,N.
\]
После выбора $\Omega_1,\ldots,\Omega_N$ положим
\[
K_0:=\overline Q\setminus\bigcup_{j=1}^{N}\Omega_j.
\]
Множество $K_0$ замкнуто и
\[
K_0\subset\widetilde Q\setminus\overline\Omega.
\]
Поэтому существует открытое $\Omega_0$ такое, что
\[
K_0\subset\Omega_0,
\qquad
\overline{\Omega_0}\subset\widetilde Q\setminus\overline\Omega.
\]
Следовательно,
\[
\overline Q\subset\bigcup_{j=0}^{N}\Omega_j,
\qquad
\overline{\Omega_0}\cap\overline\Omega=\varnothing.
\]
\end{myproof}

\begin{mytheorem}
Пусть $\Omega$ --- открытая ограниченная область,
\[
\overline\Omega\subset\bigcup_{j=1}^{N}O_j,
\qquad O_j\text{ --- открытые}.
\]
Тогда существуют функции
\[
\{\xi_j\}_{j=1}^{N},
\qquad
\xi_j\in C_0^\infty(\mathbb R^n),
\]
такие, что
\[
\operatorname{supp}\xi_j\subset O_j,
\qquad
0\le\xi_j(x)\le1,
\]
и
\[
\sum_{j=1}^{N}\xi_j(x)=1,
\qquad x\in\overline\Omega.
\]
\end{mytheorem}

\begin{myproof}
По предыдущей лемме выберем $\Omega_j$, $j=0,\ldots,N$.
Для $j=1,\ldots,N$ по первой лемме существуют функции
\[
\eta_j\in C_0^\infty(\mathbb R^n),
\qquad
\eta_j(x)\ge0,
\]
такие, что
\[
\eta_j\big|_{\overline{\Omega_j}}=1,
\qquad
\operatorname{supp}\eta_j\subset O_j.
\]
Для $j=0$ аналогично выбираем
\[
\eta_0\in C_0^\infty(\mathbb R^n),
\qquad
\eta_0(x)\ge0,
\qquad
\eta_0\big|_{\overline{\Omega_0}}=1,
\qquad
\operatorname{supp}\eta_0\subset\widetilde Q\setminus\overline\Omega.
\]

Определим, $j=1,\ldots,N$,
\[
\xi_j(x)=
\begin{cases}
\displaystyle
\frac{\eta_j(x)}{\sum\limits_{i=0}^{N}\eta_i(x)}, & x\in Q,\\[4mm]
0, & x\notin Q.
\end{cases}
\]
На $\overline Q$
\[
\sum_{i=0}^{N}\eta_i(x)\ge1.
\]
Поэтому
\[
\xi_j\in C_0^\infty(\mathbb R^n),
\qquad
\operatorname{supp}\xi_j\subset O_j.
\]
Так как при $x\in\overline\Omega$ имеем $\eta_0(x)=0$, то
\[
\sum_{j=1}^{N}\xi_j(x)=1,
\qquad x\in\overline\Omega.
\]
\end{myproof}

% ============================================================
% ГЛАВА IV
% ============================================================

\chapter{Вариационное исчисление}
В этой главе мы рассмотрим задачи вариационного исчисления, в частности, задачи поиска экстремумов функционалов.
Рассмотрим функционал
\[
J(y)=\int_a^b L\bigl(x,y(x),y'(x)\bigr)\,\dd x,
\]
где
\[
L\in C(\mathbb R^3),
\]
и задачу поиска минимума или максимума $J$ на множестве
\[
\left\{y\in C^1[a,b]:\ y(a)=A,\ y(b)=B\right\}.
\]

\section{Интегральные функционалы}

\subsection{Функционалы в нормированных пространствах}

Пусть $X$ --- линейное нормированное пространство и
\[
J:X\to\mathbb R.
\]

\begin{mydef}
Вариацией функционала $J$ в точке $f$ в направлении $h$ называется
\[
\delta J(f,h)
=
\left.\frac{\dd}{\dd\alpha}J(f+\alpha h)\right|_{\alpha=0}
\]
\end{mydef}

\begin{myremark}
\begin{enumerate}
\item Производная по Гато:
\[
\delta J(f,h)=J'_G[f]h.
\]
\item Вариация однородна по $h$.

\[
\delta J(f,\lambda h)
=
\lim_{\alpha\to 0}
\frac{
J(f+\alpha\lambda h)-J(f)
}{
\alpha
}=\begin{vmatrix}
    \beta=\alpha\lambda
\end{vmatrix}
=
\lim_{\beta\to 0}
\frac{
J(f+\beta h)-J(f)
}{
\beta/\lambda
}
=\]\[=
\lambda
\lim_{\beta\to 0}
\frac{
J(f+\beta h)-J(f)
}{
\beta
}
=
\lambda \delta J(f,h).
\]

Таким образом,
\[
\delta J(f,\lambda h)
=
\lambda\,\delta J(f,h).
\]
\end{enumerate}
\end{myremark}

\begin{myexe}
Линейности по $h$ в общем случае нет.
\end{myexe}

Пусть задан линейный функционал
\[
\varphi:X\to\mathbb R
\]

\begin{mydef}
Функционал $\varphi$ называется ограниченным, если существует $C>0$ такое, что
\[
|\varphi(x)|\le C\|x\|_X,
\qquad \forall x\in X.
\]
\end{mydef}

\begin{myexample}
Пусть
\[
X=C[-1,1],
\qquad
\|f\|_1=\int_{-1}^{1}|f(x)|\,\dd x,
\]
и
\[
\varphi(f)=f(0).
\]
Рассмотрим функции
\[
f_n(x)=
\begin{cases}
1-n|x|, & |x|\le \dfrac1n,\\[2mm]
0, & |x|>\dfrac1n.
\end{cases}
\]
Тогда
\[
\varphi(f_n)=1,
\qquad
\|f_n\|_1=\frac1n.
\]

\begin{center}
\begin{tikzpicture}[>=Stealth,scale=1]
\draw[->] (-2.2,0)--(2.2,0);
\draw[->] (0,-0.25)--(0,1.5);
\draw[thick] (-1.2,0)--(0,1)--(1.2,0);
\node[below] at (-1.2,0) {$-\frac1n$};
\node[below] at (1.2,0) {$\frac1n$};
\node[left] at (0,1) {$1$};
\end{tikzpicture}
\end{center}
\end{myexample}

\begin{myexe}
Если $\dim X<\infty$, то всякий линейный функционал на $X$ ограничен.
\end{myexe}

\begin{mydef}
Сопряжённым пространством к $X$ называется
\[
X'=X^*=
\{\text{линейные ограниченные функционалы на }X\}.
\]
Норма задаётся формулой
\[
\|\varphi\|_{X'}
=
\sup_{f\ne0}\frac{|\varphi(f)|}{\|f\|_X}.
\]
\end{mydef}

\begin{myexe}
\begin{enumerate}
\item Это действительно норма.
\item $X'$ --- банахово пространство.
\end{enumerate}
\end{myexe}



\begin{mydef}
    Пусть $J:X\to\mathbb R$ --- функционал и $f\in X$.
Говорят, что $J$ дифференцируем по Фреше в точке $f$, если существует $\varphi\in X'$ такое, что
\[
J[f+h]
=
J[f]+\varphi(h)+o(\|h\|_X),
\qquad h\to0.
\]
При этом
\[
J'_F[f]h=\varphi(h).
\]
\end{mydef}

\begin{myexe}
\begin{enumerate}
\item Если $\varphi$ существует, то оно единственно.
\item Если существует $J'_F[f]$, то
\[
J'_G[f]h=J'_F[f]h
\qquad \forall h.
\]
\item Существование $J'_G[f]h$ для всех $h$ не влечёт существования $J'_F[f]$.
\end{enumerate}
\end{myexe}

\begin{mydef}
Говорят, что
\[
J\in C^1(X),
\]
если для каждого $f\in X$ существует $J'_F[f]\in X'$ и отображение
\[
f\longmapsto J'_F[f]
\]
непрерывно как отображение $X\to X'$.
\end{mydef}

% ------------------------------------------------------------
% 03.09.26
% ------------------------------------------------------------

Пусть $X$ --- нормированное пространство и
\[
J:X\to\mathbb R.
\]

\begin{mydef}
Точка $f$ называется точкой локального минимума $J$, если
\[
\exists\delta>0:\qquad
J[f]\le J[g]
\quad\forall g\in X:\ \|g-f\|_X<\delta.
\]
Точка $f$ называется точкой строгого локального минимума, если
\[
\exists\delta>0:\qquad
J[f]<J[g]
\quad\forall g\ne f:\ \|g-f\|_X<\delta.
\]
Точка $f$ называется точкой глобального (абсолютного) минимума, если
\[
J[f]\le J[g]
\qquad\forall g\in X.
\]
\end{mydef}

\begin{mylemma}
Пусть $f$ --- локальный экстремум $J$ и вариация $\delta J(f,h)$ существует.
Тогда
\[
\delta J(f,h)=0
\qquad\forall h\in X.
\]
\end{mylemma}

\begin{myproof}
Положим
\[
\psi(\alpha)=J(f+\alpha h).
\]
Точка $\alpha=0$ является локальным экстремумом $\psi$. Если $\psi'(0)$ существует, то
\[
\psi'(0)=0.
\]
Следовательно,
\[
\delta J(f,h)=0.
\]
\end{myproof}

\begin{mydef}
Точка $f$ называется стационарной точкой функционала $J$, если
\[
\delta J(f,h)=0
\qquad\forall h.
\]
\end{mydef}

\subsection{Пример пространства $C^1$ и непрерывность интегрального функционала}

Пусть
\[
X=C^1([a,b],\mathbb R^n),
\qquad
\|f\|_X=
\max_{[a,b]}|f'(x)|+
\max_{[a,b]}|f(x)|.
\]

\begin{myexe}
\begin{enumerate}
\item $\|\cdot\|_X$ является нормой.
\item $X$ --- банахово пространство.
\end{enumerate}
\end{myexe}

\begin{mytheorem}
Пусть
\[
L\in C([a,b]\times\mathbb R^n\times\mathbb R^n),
\]
и
\[
J[y]=\int_a^bL\bigl(x,y(x),y'(x)\bigr)\,\dd x.
\]
Тогда $J:X\to\mathbb R$ непрерывен.
\end{mytheorem}

\begin{myproof}
Зафиксируем $y_0\in X$ и числа $R_0,R_1$ такие, что
\[
|y_0(x)|\le R_0,
\qquad
|y_0'(x)|\le R_1,
\qquad x\in[a,b].
\]
Функция $L$ непрерывна на компакте
\[
[a,b]\times\overline B_{R_0+1}\times\overline B_{R_1+1},
\]
поэтому равномерно непрерывна на нём.

Для каждого $\varepsilon>0$ существует $\delta\in(0,1)$ такое, что из
\[
|x-x_0|<\delta,
\qquad
|y-y_0|<\delta,
\qquad
|v-v_0|<\delta
\]
следует
\[
|L(x,y,v)-L(x_0,y_0,v_0)|<\frac{\varepsilon}{b-a}.
\]

Если
\[
\|y-y_0\|_X<\delta,
\]
то
\[
|y(x)-y_0(x)|<\delta,
\qquad
|y'(x)-y_0'(x)|<\delta,
\]
и потому
\[
\left|L(x,y(x),y'(x))-L(x,y_0(x),y_0'(x))\right|
<\frac{\varepsilon}{b-a}.
\]
Следовательно,
\[
|J[y]-J[y_0]|
\le
\int_a^b
\left|L(x,y,y')-L(x,y_0,y_0')\right|\,\dd x
<\varepsilon.
\]
\end{myproof}

\begin{mytheorem}
Пусть
\[
L\in C^1([a,b]\times\mathbb R^n\times\mathbb R^n),
\]
и
\[
J[y]=\int_a^bL\bigl(x,y(x),y'(x)\bigr)\,\dd x.
\]
Тогда
\[
J\in C^1(X),
\]
а вариация имеет вид
\[
\boxed{
\delta J(y;h)
=
\int_a^b
\left(
\left\langle\nabla_yL(x,y,y'),h(x)\right\rangle
+
\left\langle\nabla_vL(x,y,y'),h'(x)\right\rangle
\right)\dd x.}
\]
\end{mytheorem}

\begin{myproof}
Используем разложение
\[
L(x,y+h,v+w)
=
L(x,y,v)
+
\langle\nabla_yL(x,y,v),h\rangle
+
\langle\nabla_vL(x,y,v),w\rangle
+
o(|h|+|w|),
\]
при $h,w\to0$ в $\mathbb R^n$.
Подставляя
\[
y=y(x),\qquad h=h(x),\qquad v=y'(x),\qquad w=h'(x)
\]
и интегрируя по $[a,b]$, получаем
\[
J[y+h]
=
J[y]
+
\int_a^b
\left(
\langle\nabla_yL(x,y,y'),h\rangle
+
\langle\nabla_vL(x,y,y'),h'\rangle
\right)\dd x
+
o(\|h\|_X).
\]
\end{myproof}

\begin{myremark}
Отображение
\[
y\longmapsto J'_F[y],
\qquad X\to X',
\]
непрерывно.
\end{myremark}

\begin{myremark}
При $n=1$
\[
\delta J(y;h)
=
\int_a^b
\left(
\frac{\partial L(x,y,y')}{\partial y}\,h(x)
+
\frac{\partial L(x,y,y')}{\partial v}\,h'(x)
\right)\dd x.
\]
\end{myremark}

\subsection{Леммы}

\begin{mylemma}
Пусть $\Omega\subset\mathbb R^n$ --- открытое множество и
\[
f\in C(\Omega).
\]
Если
\[
\int_\Omega f(x)h(x)\,\dd x=0
\qquad
\forall h\in C_0^\infty(\Omega),
\]
то
\[
f\equiv0.
\]
\end{mylemma}

\begin{myproof}
Предположим противное. Тогда существует $a\in\Omega$ такое, что, например,
\[
f(a)>0.
\]
Тогда существует $\delta>0$ такое, что
\[
B_\delta(a)\subset\Omega,
\qquad
f(x)>0\quad\forall x\in B_\delta(a).
\]
Выберем
\[
h(x)=
\begin{cases}
\displaystyle
\exp\!\left(-\frac{1}{\delta^2-|x-a|^2}\right), & |x-a|<\delta,\\[2mm]
0, & |x-a|\ge\delta.
\end{cases}
\]
Тогда $h\in C_0^\infty(\Omega)$ и
\[
\int_\Omega f h\,\dd x
=
\int_{B_\delta(a)}fh\,\dd x>0,
\]
противоречие.
\end{myproof}

\begin{mylemma}
Пусть
\[
g\in C([a,b],\mathbb R^n).
\]
Тогда
\[
\int_a^b\langle g(x),h'(x)\rangle\,\dd x=0
\]
для всех
\[
h\in C^1([a,b],\mathbb R^n),
\qquad
h(a)=h(b)=0,
\]
тогда и только тогда, когда
\[
g(x)=\mathrm{const}.
\]
\end{mylemma}

\begin{myproof}
Обратное утверждение очевидно.
Для прямого положим
\[
c:=\frac1{b-a}\int_a^b g(x)\,\dd x
\]
и
\[
h(x)=\int_a^x(g(t)-c)\,\dd t.
\]
Тогда
\[
h'(x)=g(x)-c,
\qquad
h(a)=h(b)=0.
\]
Следовательно,
\[
0
=
\int_a^b\langle g(x),h'(x)\rangle\,\dd x
-
\int_a^b\langle c,h'(x)\rangle\,\dd x
=
\int_a^b|g(x)-c|^2\,\dd x.
\]
Отсюда
\[
g(x)=c.
\]
\end{myproof}

\begin{mylemma}
Пусть
\[
g\in C[a,b].
\]
Если
\[
\int_a^b g(x)h''(x)\,\dd x=0
\]
для всех $h\in C^2[a,b]$, удовлетворяющих
\[
h(a)=h(b)=h'(a)=h'(b)=0,
\]
то
\[
\boxed{g(x)=Ax+B.}
\]
\end{mylemma}

\begin{myproof}
Положим
\[
h(x)
=
\int_a^x\int_a^t
\bigl(g(\tau)-\alpha(\tau-a)-\beta\bigr)
\,\dd\tau\,\dd t,
\qquad
\alpha,\beta\in\mathbb R.
\]
Тогда
\[
h(a)=0,
\qquad
h'(x)=\int_a^x\bigl(g(\tau)-\alpha(\tau-a)-\beta\bigr)\,\dd\tau,
\qquad
h'(a)=0,
\]
и
\[
h''(x)=g(x)-\alpha(x-a)-\beta.
\]
Условия $h'(b)=0$ и $h(b)=0$ дают систему относительно $\alpha,\beta$:
\[
\int_a^b g(\tau)\,\dd\tau
-
\alpha\frac{(b-a)^2}{2}
-
\beta(b-a)=0,
\]
\[
\int_a^b\int_a^t g(\tau)\,\dd\tau\,\dd t
-
\alpha\frac{(b-a)^3}{6}
-
\beta\frac{(b-a)^2}{2}=0.
\]
Её определитель равен
\[
\begin{vmatrix}
\dfrac{(b-a)^2}{2} & b-a\\[2mm]
\dfrac{(b-a)^3}{6} & \dfrac{(b-a)^2}{2}
\end{vmatrix}
=
\frac{(b-a)^4}{12}>0,
\]
поэтому $\alpha,\beta$ существуют.

Тогда
\[
0
=
\int_a^b g(x)h''(x)\,\dd x
-
\int_a^b\bigl(\alpha(x-a)+\beta\bigr)h''(x)\,\dd x
\]
и, поскольку граничные значения $h,h'$ равны нулю,
\[
\int_a^b
\bigl(g(x)-\alpha(x-a)-\beta\bigr)h''(x)\,\dd x=0.
\]
Но
\[
h''(x)=g(x)-\alpha(x-a)-\beta,
\]
следовательно,
\[
g(x)=\alpha(x-a)+\beta=Ax+B.
\]
\end{myproof}

\begin{myexe}
Пусть $g\in C[a,b]$, $k\in\mathbb N$ и
\[
\int_a^b g(x)h^{(k)}(x)\,\dd x=0
\]
для всех $h\in C^k[a,b]$, удовлетворяющих
\[
h(a)=h'(a)=\cdots=h^{(k-1)}(a)=0,
\]
\[
h(b)=h'(b)=\cdots=h^{(k-1)}(b)=0.
\]
Тогда $g$ --- многочлен степени не выше $k-1$.
\end{myexe}

\subsection{Задача с фиксированными граничными точками}

\begin{mytheorem}
Пусть
\[
L\in C^1([a,b]\times\mathbb R^n\times\mathbb R^n),
\]
\[
J[y]=\int_a^bL(x,y,y')\,\dd x,
\]
и $y_0$ --- локальный экстремум $J$ на множестве
\[
\left\{y\in C^1([a,b],\mathbb R^n):\ y(a)=A,\ y(b)=B\right\}.
\]
Тогда
\[
\boxed{
\frac{\dd}{\dd x}
\nabla_vL\bigl(x,y_0(x),y_0'(x)\bigr)
=
\nabla_yL\bigl(x,y_0(x),y_0'(x)\bigr).}
\]
\end{mytheorem}

\begin{myproof}
Пусть
\[
h\in C^1([a,b],\mathbb R^n),
\qquad h(a)=h(b)=0.
\]
Тогда $\alpha=0$ --- локальный экстремум функции $J(y_0+\alpha h)$, поэтому
\[
\delta J(y_0;h)=0.
\]
Следовательно,
\[
\int_a^b
\left(
\langle\nabla_yL(x,y_0,y_0'),h\rangle
+
\langle\nabla_vL(x,y_0,y_0'),h'\rangle
\right)\dd x=0.
\]
Положим
\[
G(x):=\int_a^x\nabla_yL\bigl(t,y_0(t),y_0'(t)\bigr)\,\dd t.
\]
Тогда
\[
G\in C^1([a,b],\mathbb R^n)
\]
и
\[
\int_a^b\langle\nabla_yL,h\rangle\,\dd x
=
\langle G(x),h(x)\rangle\Big|_a^b
-
\int_a^b\langle G(x),h'(x)\rangle\,\dd x.
\]
Граничный член равен нулю, поэтому
\[
\int_a^b
\left\langle
\nabla_vL(x,y_0,y_0')-G(x),h'(x)
\right\rangle\dd x=0.
\]
По лемме Дюбуа--Реймона
\[
\nabla_vL(x,y_0,y_0')=G(x)+\mathrm{const}.
\]
Следовательно,
\[
\frac{\dd}{\dd x}\nabla_vL(x,y_0,y_0')
=
G'(x)
=
\nabla_yL(x,y_0,y_0').
\]
\end{myproof}

\begin{myremark}
Из доказательства следует, что $y_0$ обладает дополнительной гладкостью в тех случаях, когда уравнение Эйлера--Лагранжа позволяет выразить $y_0'$.
\end{myremark}

\begin{myexample}
Пусть $n=1$, $[a,b]=[-1,1]$ и
\[
J[y]=\int_{-1}^{1}y(x)^2\bigl(y'(x)-2x\bigr)^2\,\dd x,
\]
\[
y(-1)=0,
\qquad
y(1)=1.
\]
Очевидно,
\[
J[y]\ge0.
\]
Функция
\[
y_0(x)=
\begin{cases}
0, & x\le0,\\
x^2, & x>0
\end{cases}
\]
даёт
\[
J[y_0]=0.
\]
При этом
\[
y_0'(x)=
\begin{cases}
0, & x\le0,\\
2x, & x>0,
\end{cases}
\qquad
\Rightarrow\qquad y_0\notin C^2.
\]

\begin{center}
\begin{tikzpicture}[>=Stealth,scale=1]
\draw[->] (-2,0)--(2.2,0);
\draw[->] (0,-0.3)--(0,1.8);
\draw[thick] (-1.5,0)--(0,0);
\draw[thick,domain=0:1.25,smooth,variable=\x] plot ({\x},{\x*\x});
\node[below] at (-1.5,0) {$-1$};
\node[below] at (1,0) {$1$};
\end{tikzpicture}
\end{center}
\end{myexample}

\begin{mytheorem}
Пусть выполнены условия предыдущей теоремы,
\[
L\in C^2([a,b]\times\mathbb R^n\times\mathbb R^n),
\]
и
\[
\det\left(
\frac{\partial^2L(x,y,v)}{\partial v_j\partial v_k}
\right)_{j,k=1}^{n}
\ne0
\qquad \forall x,y,v.
\]
Тогда
\[
y_0\in C^2([a,b],\mathbb R^n).
\]
\end{mytheorem}

\begin{myproof}
Рассмотрим отображение
\[
\Phi:(x,y,v)\longmapsto
\bigl(x,y,\nabla_vL(x,y,v)\bigr).
\]
Имеем
\[
\Phi:[a,b]\times\mathbb R^{2n}	o[a,b]\times\mathbb R^{2n},
\qquad
\Phi\in C^1.
\]
Матрица $D\Phi(x,y,v)$ имеет блочно-треугольный вид, поэтому
\[
\det D\Phi
=
\det\left(
\frac{\partial^2L}{\partial v_j\partial v_k}
\right)
e0.
\]

Для любого $x_0\in(a,b)$ по теореме об обратной функции существует локальное обратное отображение
\[
\Psi\in C^1,
\]
такое, что
\[
\Psi\bigl(x,y,\nabla_vL(x,y,v)\bigr)=(x,y,v).
\]
Следовательно, в окрестности $x_0$
\[
y_0'(x)
=
\Psi_{n+1:\,2n}
\left(
 x,y_0(x),
 \nabla_vL(x,y_0(x),y_0'(x))
\right).
\]
Правая часть принадлежит $C^1$, поэтому
\[
y_0'\in C^1,
\qquad
y_0\in C^2.
\]
\end{myproof}

\begin{myexe}
Разобраться с концами $x_0=a,b$.
\end{myexe}

\subsection{Задача со свободными концами}

\begin{mytheorem}
Пусть
\[
L\in C^1([a,b]\times\mathbb R^n\times\mathbb R^n),
\]
и $y_0$ --- локальный экстремум функционала
\[
J[y]=\int_a^bL(x,y,y')\,\dd x
\]
на всём пространстве $X=C^1([a,b],\mathbb R^n)$.
Тогда
\[
\boxed{
\frac{\dd}{\dd x}\nabla_vL(x,y_0,y_0')
=
\nabla_yL(x,y_0,y_0')}
\]
и
\[
\boxed{
\nabla_vL(a,y_0(a),y_0'(a))
=
\nabla_vL(b,y_0(b),y_0'(b))
=0.}
\]
\end{mytheorem}

\begin{myproof}
Для любого $h\in X$
\[
\delta J(y_0,h)=0,
\]
то есть
\[
\int_a^b
\left(
\langle\nabla_yL,h\rangle
+
\langle\nabla_vL,h'\rangle
\right)\dd x=0.
\]
Если сначала взять $h(a)=h(b)=0$, то получаем уравнение Эйлера--Лагранжа.
После интегрирования по частям с учётом этого уравнения остаётся
\[
\left.\langle\nabla_vL(x,y_0,y_0'),h(x)\rangle\right|_a^b=0.
\]
Выбирая $h$ с произвольными значениями на концах, получаем естественные граничные условия
\[
\nabla_vL(a,y_0(a),y_0'(a))=0,
\qquad
\nabla_vL(b,y_0(b),y_0'(b))=0.
\]
\end{myproof}

\begin{myremark}
Уравнение Эйлера--Лагранжа вместе с естественными граничными условиями эквивалентно стационарности $y_0$.
\end{myremark}

\begin{mytheorem}
Если дополнительно
\[
L\in C^2([a,b]\times\mathbb R^{2n}),
\qquad
\det\left(
\frac{\partial^2L}{\partial v_j\partial v_k}
\right)\ne0,
\]
то
\[
y_0\in C^2([a,b],\mathbb R^n).
\]
\end{mytheorem}

\begin{myremark}
Уравнение
\[
\frac{\dd}{\dd x}\nabla_vL(x,y_0,y_0')
=
\nabla_yL(x,y_0,y_0')
\]
называется уравнением Эйлера--Лагранжа, а
\[
\nabla_vL(x,y_0,y_0')\big|_{x=a,b}=0
\]
--- естественными граничными условиями.
Для $n$ неизвестных функций система второго порядка содержит $2n$ граничных условий.
\end{myremark}

\subsection{Частные случаи уравнения Эйлера--Лагранжа}

\begin{enumerate}
\item Если
\[
L=L(x,y),
\]
то
\[
\nabla_yL(x,y_0(x))=0.
\]
Это не дифференциальное уравнение.

\item Если $L$ не зависит от $y_j$, то
\[
\frac{\dd}{\dd x}
\frac{\partial L}{\partial v_j}(x,y_0,y_0')=0,
\]
то есть
\[
\frac{\partial L}{\partial v_j}(x,y_0,y_0')=\mathrm{const}.
\]
В соответствующем случае
\[
\nabla_vL(x,y_0,y_0')=\vec c,
\]
откуда
\[
y_0'(x)=f(x,\vec c),
\qquad
 y_0(x)=\int_a^x f(t,\vec c)\,\dd t+\vec c_2.
\]

\item Если
\[
L=L(y,v),
\]
то
\[
\boxed{
\frac{\dd}{\dd x}
\left(
L(y_0,y_0')-
\langle y_0',\nabla_vL(y_0,y_0')\rangle
\right)=0.}
\]
\end{enumerate}

\begin{myexample}
При $n=1$ рассмотрим
\[
\int \frac{\sqrt{1+y'^2}}{y}\,\dd x,
\qquad
L(y,v)=\frac{\sqrt{1+v^2}}{y},
\qquad y>0.
\]
Имеем
\[
\frac{\partial L}{\partial v}
=
\frac{v}{y\sqrt{1+v^2}}.
\]
Следовательно,
\[
L-y'\frac{\partial L}{\partial v}
=
\frac{\sqrt{1+y'^2}}{y}
-
\frac{y'^2}{y\sqrt{1+y'^2}}
=
\frac1{y\sqrt{1+y'^2}}.
\]
Пусть
\[
\frac1{y\sqrt{1+y'^2}}=\frac1R,
\qquad R>0.
\]
Тогда
\[
y^2(1+y'^2)=R^2,
\]
\[
y'^2=\frac{R^2}{y^2}-1,
\]
\[
y'=\pm\frac{\sqrt{R^2-y^2}}{y}.
\]
Отсюда
\[
\int\frac{y\,\dd y}{\sqrt{R^2-y^2}}
=
\pm\int\dd x,
\]
и
\[
\sqrt{R^2-y^2}=\pm(x-a).
\]
Следовательно,
\[
\boxed{(x-a)^2+y^2=R^2.}
\]
\end{myexample}

% ============================================================
% 10.09.26
% ============================================================

\section{Условный экстремум}

\subsection{Касательные к гиперповерхности}

Пусть $X$ --- линейное нормированное пространство,
\[
G\in C^1(X),
\qquad
M=\{f\in X:\ G[f]=0\}.
\]
Пусть
\[
G'_F[f]\ne0
\qquad \forall f\in M.
\]
Тогда $M$ --- гиперповерхность.

Касательное пространство к $M$ в точке $u\in M$ имеет вид
\[
\{u+h:\ G'_F[u]h=0\}
=
 u+\ker G'_F[u].
\]

В случае $X=\mathbb R^n$
\[
G[u+h]
=
G[u]+\langle\nabla G(u),h\rangle+o(|h|),
\qquad h\to0.
\]

\begin{mytheorem}
Пусть
\[
G\in C^1(X),
\qquad
M=\{f:\ G[f]=0\},
\]
и $u\in M$ не является стационарной точкой $G$.
Тогда условие
\[
\delta G(u;h)=0
\]
эквивалентно существованию $\varepsilon>0$ и кривой
\[
w:(-\varepsilon,\varepsilon)\to X,
\qquad w\in C^1,
\]
такой, что
\[
w(t)\in M,
\qquad
w(0)=u,
\qquad
w'(0)=h.
\]
\end{mytheorem}

\begin{myproof}
Так как $u$ не является стационарной точкой $G$, существует $v\in X$ такое, что
\[
\delta G(u;v)\ne0.
\]
Положим
\[
\psi(s,t)=G[u+sv+th].
\]
Тогда
\[
\psi(0,0)=G[u]=0.
\]
При $t=0$
\[
\psi(s,0)=G[u+sv]
=
G[u]+s\,\delta G(u;v)+o(s),
\]
поэтому
\[
\frac{\partial\psi(0,0)}{\partial s}
=
\delta G(u;v)\ne0.
\]
А
\[
\psi(0,t)=G[u+th]
=
G[u]+t\,\delta G(u;h)+o(t),
\]
поэтому
\[
\frac{\partial\psi(0,0)}{\partial t}
=
\delta G(u;h)=0.
\]

По теореме о неявной функции существует функция $s=s(t)$,
\[
\psi(s(t),t)=0,
\qquad
s(0)=0,
\]
и
\[
s'(0)
=
-\frac{\partial_t\psi}{\partial_s\psi}(0,0)
=0.
\]
Следовательно,
\[
s(t)=o(t),
\qquad t\to0.
\]
Положим
\[
w(t)=u+s(t)v+th.
\]
Тогда
\[
G[w(t)]=0,
\qquad
w(t)\in M,
\qquad
w(0)=u,
\]
и
\[
w'(0)
=
\lim_{t\to0}\frac{w(t)-w(0)}{t}
=
\lim_{t\to0}\left(\frac{s(t)}t v+h\right)
=h.
\]
Обратное утверждение получается дифференцированием $G[w(t)]=0$ при $t=0$.
\end{myproof}

\begin{mylemma}
Пусть $J,G\in C^1(X)$,
\[
M=\{f\in X:\ G[f]=0\}.
\]
Пусть $u$ --- локальный экстремум $J$ на $M$, причём $u$ не является стационарной точкой $G$.
Если
\[
\delta G(u;h)=0,
\]
то
\[
\delta J(u;h)=0.
\]
\end{mylemma}

\begin{myproof}
По предыдущей теореме существует кривая
\[
w(t)=u+s(t)v+th,
\qquad
w(t)\in M,
\qquad
w(0)=u,
\qquad
w'(0)=h.
\]
Поскольку $u$ --- локальный экстремум $J$ на $M$, точка $t=0$ является локальным экстремумом функции $J[w(t)]$, поэтому
\[
\left.\frac{\dd}{\dd t}J[w(t)]\right|_{t=0}=0.
\]
Разложение
\[
J[w(t)]
=
J[u]
+s(t)\delta J(u;v)
+t\delta J(u;h)
+o(\|s(t)v+th\|_X)
\]
и $s(t)=o(t)$ дают
\[
\delta J(u;h)=0.
\]
\end{myproof}

\begin{mytheorem}
Пусть $J,G\in C^1(X)$,
\[
M=\{f:\ G[f]=0\},
\]
и $u$ --- локальный экстремум $J$ на $M$, причём $u$ не является стационарной точкой $G$.
Тогда существует $\lambda\in\mathbb R$ такое, что $u$ является стационарной точкой функционала
\[
J+\lambda G.
\]
\end{mytheorem}

\begin{myproof}
Выберем $v\in X$ так, что
\[
\delta G(u;v)\ne0,
\]
и положим
\[
\lambda:=-\frac{\delta J(u;v)}{\delta G(u;v)}.
\]
Для произвольного $y\in X$ положим
\[
h:=y-\frac{\delta G(u;y)}{\delta G(u;v)}v.
\]
Тогда
\[
\delta G(u;h)=0.
\]
По предыдущей лемме
\[
\delta J(u;h)=0.
\]
Следовательно,
\[
\delta J(u;y)
-
\frac{\delta G(u;y)}{\delta G(u;v)}\delta J(u;v)=0,
\]
то есть
\[
\delta(J+\lambda G)(u;y)=0
\qquad\forall y\in X.
\]
\end{myproof}

\begin{mytheorem}
Пусть
\[
J,G_1,G_2,\ldots,G_m\in C^1(X),
\]
\[
M=\{f\in X:\ G_1[f]=\cdots=G_m[f]=0\},
\]
и $u$ --- локальный экстремум $J$ на $M$.
Рассмотрим линейный оператор
\[
\vec G'_F[u]:X\to\mathbb R^m,
\qquad
h\longmapsto
\begin{pmatrix}
\delta G_1(u;h)\\
\vdots\\
\delta G_m(u;h)
\end{pmatrix}.
\]
Если
\[
\operatorname{Im}\vec G'_F[u]=\mathbb R^m,
\]
то существуют
\[
\lambda_1,\ldots,\lambda_m\in\mathbb R
\]
такие, что $u$ является стационарной точкой функционала
\[
J+\lambda_1G_1+\cdots+\lambda_mG_m.
\]
\end{mytheorem}

\begin{myremark}
Числа $\lambda_1,\ldots,\lambda_m$ называются множителями Лагранжа.
\end{myremark}

\begin{myexe}
\begin{enumerate}
\item При условиях предыдущей теоремы для всякого касательного вектора $h$, то есть
\[
\delta G_j(u;h)=0,
\qquad j=1,\ldots,m,
\]
существует кривая $w\in C^1(( -\varepsilon,\varepsilon);X)$ такая, что
\[
w(t)\in M,
\qquad
w(0)=u,
\qquad
w'(0)=h.
\]
\item Если $h$ --- касательный вектор, то
\[
\delta J(u;h)=0.
\]
\item Доказать предыдущую теорему.
\end{enumerate}
\end{myexe}

\begin{myexample}
Пусть $\Omega$ --- область с границей
\[
\partial\Omega=\Gamma,
\]
на которой распределён заряд с плотностью $\rho(x)$,
\[
\rho\in C(\Gamma),
\qquad
\int_\Gamma\rho(x)\,\dd S(x)=q.
\]
Потенциальная энергия имеет вид
\[
J[\rho]
=
\iint_{\Gamma\times\Gamma}
\frac{\rho(x)\rho(y)}{|x-y|}
\,\dd S(x)\dd S(y),
\]
и требуется найти её минимум при условии
\[
G[\rho]
=
\int_\Gamma\rho(x)\,\dd S(x)-q=0.
\]
По теореме о множителе Лагранжа существует $\lambda\in\mathbb R$ такое, что
\[
\delta(J+\lambda G)(\rho;h)=0
\qquad\forall h.
\]
Имеем
\[
\left.
\frac{\dd}{\dd\alpha}
\left[
\iint_{\Gamma\times\Gamma}
\frac{(\rho(x)+\alpha h(x))(\rho(y)+\alpha h(y))}{|x-y|}
\,\dd S(x)\dd S(y)
+
\lambda\left(\int_\Gamma(\rho+\alpha h)\,\dd S-q\right)
\right]
\right|_{\alpha=0}
=0.
\]
Отсюда
\[
\iint_{\Gamma\times\Gamma}
\frac{\rho(x)h(y)+\rho(y)h(x)}{|x-y|}
\,\dd S(x)\dd S(y)
+
\lambda\int_\Gamma h(x)\,\dd S(x)=0,
\]
то есть
\[
\int_\Gamma
\left(
2\int_\Gamma\frac{\rho(y)}{|x-y|}\,\dd S(y)+\lambda
\right)
 h(x)\,\dd S(x)=0.
\]
Следовательно,
\[
\boxed{
\int_\Gamma\frac{\rho(y)}{|x-y|}\,\dd S(y)
=-\frac{\lambda}{2},
\qquad x\in\Gamma.}
\]
Это задача Робена.
Если
\[
\varphi(x)=\int_\Gamma\frac{\rho(y)}{|x-y|}\,\dd S(y),
\]
то
\[
\Delta\varphi(x)=0\quad\text{в }\Omega,
\qquad
\varphi\big|_\Gamma=-\frac\lambda2,
\]
откуда
\[
\varphi(x)=-\frac\lambda2
\qquad\text{в }\Omega.
\]
\end{myexample}

% ------------------------------------------------------------
% 15.09.26
% ------------------------------------------------------------

\subsection{Изопериметрическая задача}

Пусть
\[
X=C^1([a,b],\mathbb R^n).
\]

\begin{mytheorem}
Пусть
\[
L,g_1,\ldots,g_m
\in
C^1([a,b]\times\mathbb R^n\times\mathbb R^n),
\]
\[
G_j[y]
=
\int_a^b g_j(x,y(x),y'(x))\,\dd x,
\qquad
J[y]
=
\int_a^bL(x,y,y')\,\dd x,
\]
\[
M=\{y\in X:\ G_j[y]=0,\ j=1,\ldots,m\}.
\]
Пусть $y_0$ --- локальный экстремум $J$ на $M$ и
\[
\operatorname{Im}\vec G'_F[y_0]=\mathbb R^m.
\]
Тогда существуют
\[
\lambda_1,\ldots,\lambda_m\in\mathbb R
\]
такие, что для
\[
F(x,y,v)
=
L(x,y,v)+\sum_{j=1}^{m}\lambda_jg_j(x,y,v)
\]
выполняются условия
\[
\boxed{
\frac{\dd}{\dd x}\nabla_vF(x,y_0,y_0')
=
\nabla_yF(x,y_0,y_0')}
\]
и, в случае свободных концов,
\[
\boxed{
\nabla_vF(x,y_0(x),y_0'(x))\big|_{x=a,b}=0.}
\]
\end{mytheorem}

\begin{myremark}
Если $y_0$ --- локальный экстремум $J$ на множестве
\[
M=
\left\{
 y\in X:
 G_j[y]=0,\ j=1,\ldots,m;
 y(a)=A,\ y(b)=B
\right\},
\]
то
\[
\frac{\dd}{\dd x}\nabla_vF=\nabla_yF.
\]
\end{myremark}

\begin{myremark}
Получается $n$ уравнений второго порядка, то есть $2n+m$ неизвестных постоянных. На них имеются $2n$ граничных условий и $m$ интегральных условий
\[
\int_a^b g_j(x,y_0,y_0')\,\dd x=0.
\]
\end{myremark}

\subsection{Задача Лагранжа}

Рассмотрим
\[
J[y]=\int_a^bL(x,y,y')\,\dd x
\]
при связях
\[
H_i(x,y(x),y'(x))=0,
\qquad i=1,\ldots,m.
\]

\begin{mytheorem}
Пусть
\[
L,H_1,\ldots,H_m
\in
C^2([a,b]\times\mathbb R^n\times\mathbb R^n),
\]
и $y_0$ --- локальный экстремум $J$ на множестве
\[
\left\{
 y\in X:
 H_i(x,y(x),y'(x))=0,
\ i=1,\ldots,m,
\ x\in[a,b];
\ y(a)=A,\ y(b)=B
\right\}.
\]
Пусть $n>m$. Предположим, что выполнено одно из условий:

\begin{enumerate}
\item $H=H(x,y,v)$ и
\[
\operatorname{rank}
\left(
\frac{\partial H_i}{\partial v_j}(x,y,v)
\right)=m
\qquad\forall x,y,v;
\]
\item $H=H(x,y)$ и
\[
\operatorname{rank}
\left(
\frac{\partial H_i}{\partial y_j}(x,y)
\right)=m
\qquad\forall x,y.
\]
\end{enumerate}

Тогда существуют функции
\[
\lambda_1,\ldots,\lambda_m\in C[a,b]
\]
такие, что при
\[
F(x,y,v)
=
L(x,y,v)+\sum_{i=1}^{m}\lambda_i(x)H_i(x,y,v)
\]
выполняется
\[
\boxed{
\frac{\dd}{\dd x}\nabla_vF(x,y_0,y_0')
=
\nabla_yF(x,y_0,y_0').}
\]
\end{mytheorem}

% ============================================================
% §3. Задача Штурма--Лиувилля
% ============================================================

\section{Задача Штурма--Лиувилля}

\subsection{Минимизирующая последовательность}

Рассмотрим функционал
\[
J[y]
=
\int_a^b
\left(
 p(x)y'(x)^2+q(x)y(x)^2
\right)\dd x,
\]
где
\[
p\in C^1[a,b],
\qquad
p(x)>0\quad\forall x,
\qquad
q\in C[a,b].
\]
Пусть
\[
M_1=
\left\{
 y\in C^1[a,b]:
 y(a)=y(b)=0,
 \int_a^b y(x)^2\,\dd x=1
\right\}.
\]

\begin{myremark}
Если $y_0$ --- локальный экстремум, то по теореме о множителе Лагранжа существует $\lambda\in\mathbb R$ такое, что $y_0$ --- стационарная точка функционала
\[
\int_a^bF\,\dd x,
\qquad
F(x,y,v)=p(x)v^2+q(x)y^2-\lambda y^2.
\]
Поэтому
\[
\frac{\dd}{\dd x}\bigl(2p(x)y_0'(x)\bigr)
=
2(q(x)-\lambda)y_0(x),
\]
то есть
\[
\boxed{
(p y_0')'+(\lambda-q)y_0=0,}
\qquad
\boxed{y_0(a)=y_0(b)=0.}
\]
\end{myremark}

Положим
\[
\lambda_1:=\inf_{M_1}J[y].
\]

\begin{mylemma}
Если
\[
q_0:=\min_{[a,b]}q(x),
\]
то
\[
\lambda_1\ge q_0.
\]
\end{mylemma}

\begin{myproof}
Для $y\in M_1$
\[
J[y]
\ge
\int_a^bq(x)y(x)^2\,\dd x
\ge
q_0\int_a^b y(x)^2\,\dd x
=q_0.
\]
\end{myproof}

\begin{mylemma}
Если
\[
f\in C^1[a,b],
\qquad
f(a)=f(b)=0,
\]
то
\[
\boxed{
J[f]\ge\lambda_1\int_a^b f(x)^2\,\dd x.}
\]
\end{mylemma}

\begin{myproof}
Положим
\[
y(x)=
\frac{f(x)}{
\sqrt{\displaystyle\int_a^b f(x)^2\,\dd x}}.
\]
Тогда $y\in M_1$, следовательно
\[
J[y]\ge\lambda_1.
\]
После умножения на $\int_a^b f^2\,\dd x$ получаем требуемое.
\end{myproof}

Существует минимизирующая последовательность
\[
\{f_m\}\subset M_1,
\qquad
J[f_m]\to\lambda_1.
\]

\begin{mytheorem}
\textbf{Теорема Арцела--Асколи.}\par
Пусть
\[
\{h_k\}_{k=1}^{\infty}\subset C[a,b]
\]
--- равномерно ограниченное и равностепенно непрерывное семейство. Тогда существует подпоследовательность
\[
\{h_{k_j}\}_{j=1}^{\infty}
\]
и функция $h_*\in C[a,b]$ такие, что
\[
h_{k_j}\to h_*
\]
равномерно на $[a,b]$.
\end{mytheorem}

\begin{mylemma}
Пусть
\[
f_m\in C^1[a,b],
\qquad
f_m(a)=0,
\]
и существует $K>0$ такое, что
\[
\int_a^b f_m'(x)^2\,\dd x\le K.
\]
Тогда существует подпоследовательность $\{f_{m_j}\}$, сходящаяся равномерно на $[a,b]$ к некоторой функции $y\in C[a,b]$.
\end{mylemma}

\begin{myproof}
Пусть $x_1<x_2$. Тогда
\[
|f_m(x_2)-f_m(x_1)|
=
\left|\int_{x_1}^{x_2}f_m'(t)\,\dd t\right|
\le
\sqrt{x_2-x_1}
\sqrt{\int_{x_1}^{x_2}|f_m'(t)|^2\,\dd t}
\le
\sqrt{x_2-x_1}\,\sqrt K.
\]
Следовательно, семейство $\{f_m\}$ равностепенно непрерывно.
Кроме того, при $x_1=a$, $x_2=x$ получаем
\[
|f_m(x)|\le\sqrt{(b-a)K},
\]
то есть семейство равномерно ограничено. По теореме Арцела--Асколи существует равномерно сходящаяся подпоследовательность.
\end{myproof}

\begin{mylemma}
Пусть
\[
p_0:=\min_{[a,b]}p(x)>0.
\]
Тогда для $f\in M_1$
\[
\boxed{
\int_a^b f'(x)^2\,\dd x
\le
\frac{J[f]-q_0}{p_0}.}
\]
\end{mylemma}

\begin{myproof}
Имеем
\[
J[f]
\ge
p_0\int_a^b f'(x)^2\,\dd x
+
q_0\int_a^b f(x)^2\,\dd x.
\]
Так как $f\in M_1$, последний интеграл равен $1$.
\end{myproof}

% ------------------------------------------------------------
% 17.09.26
% ------------------------------------------------------------

Пусть
\[
J[y]=\int_a^b(py'^2+qy^2)\,\dd x,
\qquad
p\in C^1,
\quad p>0,
\quad q\in C,
\]
и
\[
M_1=
\left\{
 y\in C^1[a,b]:
 y(a)=y(b)=0,
 \int_a^b y(x)^2\,\dd x=1
\right\}.
\]
Пусть
\[
\lambda_1:=\inf_{y\in M_1}J[y],
\qquad
J[f_m]\to\lambda_1.
\]
По предыдущим леммам существует сходящаяся подпоследовательность.

\begin{mytheorem}
Существует подпоследовательность минимизирующей последовательности $\{f_m\}\subset M_1$ и функция $y_1\in C[a,b]$ такие, что
\[
f_m(x)\to y_1(x)
\]
равномерно на $[a,b]$.
\end{mytheorem}

\subsection{Свойства предельной функции}

Имеем
\[
y_1(a)=y_1(b)=0,
\qquad
\int_a^b y_1(x)^2\,\dd x=1.
\]

Введём оператор
\[
\mathcal L
=
-\frac{\dd}{\dd x}
\left(p(x)\frac{\dd}{\dd x}\right)
+q(x),
\]
\[
\operatorname{Dom}\mathcal L
=
\{f\in C^2[a,b]:\ f(a)=f(b)=0\}.
\]

\begin{mylemma}
Для любого $h\in\operatorname{Dom}\mathcal L$
\[
\boxed{
\int_a^b y_1\,\mathcal L h\,\dd x
=
\lambda_1\int_a^b y_1h\,\dd x.}
\]
\end{mylemma}

\begin{myproof}
Пусть $h\in\operatorname{Dom}\mathcal L$ и $\tau\in\mathbb R$. По предыдущей лемме
\[
J[f_m+\tau h]
\ge
\lambda_1
\int_a^b(f_m+\tau h)^2\,\dd x.
\]
Раскрывая обе части,
\[
J[f_m]
+2\tau\int_a^b(pf_m'h'+qf_mh)\,\dd x
+\tau^2J[h]
\]
\[
\ge
\lambda_1
\int_a^b f_m^2\,\dd x
+2\tau\lambda_1\int_a^b f_mh\,\dd x
+\tau^2\lambda_1\int_a^b h^2\,\dd x.
\]
Переходя к пределу по подпоследовательности и используя
\[
\int_a^b(pf_m'h'+qf_mh)\,\dd x
=
\int_a^b f_m\,\mathcal Lh\,\dd x,
\]
получаем
\[
\lambda_1
+2\tau\int_a^b y_1\mathcal Lh\,\dd x
+\tau^2J[h]
\ge
\lambda_1
+2\tau\lambda_1\int_a^b y_1h\,\dd x
+\tau^2\lambda_1\int_a^b h^2\,\dd x.
\]
Это неравенство имеет вид
\[
A\tau^2+B\tau\ge0
\qquad\forall\tau\in\mathbb R,
\]
откуда
\[
B=0.
\]
Следовательно,
\[
\int_a^b y_1\mathcal Lh\,\dd x
=
\lambda_1\int_a^b y_1h\,\dd x.
\]
\end{myproof}

\begin{mytheorem}
Пусть $y\in C[a,b]$. Тогда
\[
y\in\operatorname{Dom}\mathcal L,
\qquad
\mathcal Ly=\lambda y
\]
тогда и только тогда, когда
\[
\int_a^b y\mathcal Lh\,\dd x
=
\lambda\int_a^b yh\,\dd x
\qquad\forall h\in\operatorname{Dom}\mathcal L.
\]
\end{mytheorem}

\begin{myproof}
Если $y\in\operatorname{Dom}\mathcal L$ и $\mathcal Ly=\lambda y$, то по симметрии
\[
(y,\mathcal Lh)
=
(\mathcal Ly,h)
=
\lambda(y,h).
\]

Обратно, положим
\[
\widetilde q(x)=q(x)-\lambda.
\]
Тогда
\[
\int_a^b y\left[-(ph')'+\widetilde qh\right]\dd x=0
\qquad\forall h\in\operatorname{Dom}\mathcal L.
\]
Выберем функции $g_1,g_2$ так, что
\[
g_1'=yp',
\qquad
 g_1\in C^1,
\]
\[
g_2''=y\widetilde q,
\qquad
 g_2\in C^2.
\]
Тогда после преобразования
\[
\int_a^b(-py+g_1+g_2)h''(x)\,\dd x=0
\]
для всех $h\in C^2[a,b]$ с
\[
h(a)=h'(a)=h(b)=h'(b)=0.
\]
По обобщённой лемме Дюбуа--Реймона
\[
-py+g_1+g_2=Ax+B.
\]
Следовательно,
\[
py=g_1+g_2-Ax-B\in C^1[a,b],
\]
откуда
\[
y\in C^1[a,b].
\]
Дифференцируя,
\[
p'y+py'=g_1'+g_2'-A.
\]
Так как $g_1'=yp'$, получаем
\[
py'=g_2'-A\in C^1[a,b],
\]
и потому
\[
y\in C^2[a,b].
\]
Далее
\[
(py')'=g_2''=\widetilde qy,
\]
то есть
\[
\mathcal Ly=\lambda y.
\]
Наконец, интегрирование по частям в исходном слабом равенстве даёт
\[
-py h'\Big|_a^b=0
\]
для всех допустимых $h$, следовательно
\[
py\Big|_a^b=0.
\]
Так как $p>0$,
\[
y(a)=y(b)=0.
\]
\end{myproof}

\begin{myexe}
Пусть $y\in C[a,b]$, $f\in C[a,b]$ и
\[
\int_a^b y\mathcal Lh\,\dd x
=
\int_a^b fh\,\dd x
\qquad\forall h\in\operatorname{Dom}\mathcal L.
\]
Тогда
\[
y\in\operatorname{Dom}\mathcal L,
\qquad
\mathcal Ly=f.
\]
\end{myexe}

\begin{mytheorem}
Пусть
\[
p\in C^1[a,b],\quad p>0,
\qquad q\in C[a,b].
\]
Тогда
\[
\inf_{M_1}J
\]
достигается: существует
\[
y_1\in M_1\cap\operatorname{Dom}\mathcal L
\]
такое, что
\[
J[y_1]=\lambda_1=\min_{M_1}J,
\qquad
\mathcal Ly_1=\lambda_1y_1.
\]
\end{mytheorem}

\begin{myproof}
По слабому равенству
\[
\mathcal Ly_1=\lambda_1y_1.
\]
Тогда
\[
J[y_1]
=
\int_a^b(py_1'^2+qy_1^2)\,\dd x
=
\int_a^b\bigl(-(py_1')'+qy_1\bigr)y_1\,\dd x
=
\lambda_1\int_a^b y_1(x)^2\,\dd x
=
\lambda_1.
\]
\end{myproof}

\subsection{Остальные собственные числа}

Положим
\[
M_2=
\left\{
 y\in M_1:
 \int_a^b yy_1\,\dd x=0
\right\},
\qquad
\lambda_2:=\inf_{M_2}J.
\]

Пусть уже построены $y_1,\ldots,y_{n-1}$. Определим
\[
M_n=
\left\{
 y\in C^1[a,b]:
 y(a)=y(b)=0,
 \int_a^b y(x)^2\,\dd x=1,
 \int_a^b y(x)y_k(x)\,\dd x=0,
\ k=1,\ldots,n-1
\right\}
\]
и
\[
\lambda_n:=\inf_{M_n}J[y].
\]
Существует минимизирующая последовательность
\[
\{f_m\}\subset M_n,
\qquad
J[f_m]\to\lambda_n.
\]
Из оценки производных снова следует существование сходящейся подпоследовательности
\[
f_{m_k}\to y_n
\]
равномерно, причём
\[
y_n\in C[a,b],
\qquad
y_n(a)=y_n(b)=0,
\]
\[
\int_a^b y_n(x)^2\,\dd x=1,
\qquad
\int_a^b y_ny_k\,\dd x=0,
\quad k=1,\ldots,n-1.
\]

\begin{mylemma}
Для любого $h\in\operatorname{Dom}\mathcal L$ существует разложение
\[
h=h_0+h_1,
\qquad
h_0,h_1\in\operatorname{Dom}\mathcal L,
\]
где
\[
h_0(x)=\sum_{k=1}^{n-1}c_ky_k(x),
\]
и
\[
\int_a^b h_1(x)y_k(x)\,\dd x=0,
\qquad k=1,\ldots,n-1.
\]
\end{mylemma}

\begin{myproof}
Для $h\in\operatorname{Dom}\mathcal L$ положим
\[
c_k:=\int_a^b h(x)y_k(x)\,\dd x,
\qquad k=1,\ldots,n-1.
\]
Затем
\[
h_0:=\sum_{k=1}^{n-1}c_ky_k,
\qquad
h_1:=h-h_0.
\]
Тогда
\[
\int_a^b h_1y_k\,\dd x
=
\int_a^b hy_k\,\dd x
-
\int_a^b h_0y_k\,\dd x
=c_k-c_k=0.
\]
\end{myproof}

\begin{mylemma}
Для любого $h\in\operatorname{Dom}\mathcal L$
\[
\boxed{
\int_a^b y_n\mathcal Lh\,\dd x
=
\lambda_n\int_a^b y_nh\,\dd x.}
\]
\end{mylemma}

\begin{myproof}
Пусть сначала $h_1\in\operatorname{Dom}\mathcal L$ удовлетворяет
\[
\int_a^b h_1y_k\,\dd x=0,
\qquad k=1,\ldots,n-1.
\]
Тогда, варьируя минимизирующую последовательность в направлении $h_1$, получаем
\[
\int_a^b y_n\mathcal Lh_1\,\dd x
=
\lambda_n\int_a^b y_nh_1\,\dd x.
\]

Для
\[
h_0=\sum_{k=1}^{n-1}c_ky_k
\]
имеем
\[
\int_a^b y_nh_0\,\dd x=0,
\]
а также
\[
\int_a^b y_n\mathcal Lh_0\,\dd x
=
\int_a^b y_n
\sum_{k=1}^{n-1}c_k\lambda_ky_k\,\dd x
=0.
\]
Складывая два равенства, получаем утверждение для произвольного $h$.
\end{myproof}

\begin{mytheorem}
Имеем
\[
y_n\in\operatorname{Dom}\mathcal L,
\qquad
\mathcal Ly_n=\lambda_ny_n,
\qquad
J[y_n]=\lambda_n,
\qquad
 y_n\in M_n.
\]
\end{mytheorem}

\begin{myproof}
Предыдущая лемма и теорема о слабом решении дают
\[
y_n\in\operatorname{Dom}\mathcal L,
\qquad
\mathcal Ly_n=\lambda_ny_n.
\]
Далее
\[
J[y_n]
=(\mathcal Ly_n,y_n)
=\lambda_n.
\]
\end{myproof}

% ------------------------------------------------------------
% 22.09.26
% ------------------------------------------------------------

\begin{mylemma}
Собственные числа удовлетворяют
\[
\boxed{
\lambda_{n+1}>\lambda_n,
\qquad
\lambda_n\to+\infty.}
\]
\end{mylemma}

\begin{myproof}
Из определения $M_{n+1}\subset M_n$ следует
\[
\lambda_{n+1}\ge\lambda_n.
\]
Если бы
\[
\lambda_{n+1}=\lambda_n,
\]
то $y_n$ и $y_{n+1}$ отвечали бы одному собственному значению, что противоречит их ортогональности.

Предположим теперь, что существует $\Lambda$ такое, что
\[
\lambda_n\le\Lambda
\qquad\forall n.
\]
Тогда
\[
J[y_n]=\lambda_n\le\Lambda,
\]
поэтому последовательность $\{y_n\}$ имеет равномерно сходящуюся подпоследовательность
\[
y_{n_k}\to z
\]
на $[a,b]$.
Следовательно,
\[
\int_a^b|y_{n_k}-y_{n_j}|^2\,\dd x\to0
\qquad k,j\to\infty.
\]
Но при $k\ne j$ из ортонормированности
\[
\int_a^b|y_{n_k}-y_{n_j}|^2\,\dd x=2,
\]
противоречие.
\end{myproof}

\begin{myremark}
При $n\to\infty$
\[
\lambda_n=Cn^2+O(1),
\]
где
\[
\boxed{
C=
\left(
\frac{\pi}{\displaystyle\int_a^b\frac{\dd x}{\sqrt{p(x)}}}
\right)^2.}
\]
\end{myremark}

\subsection{Полнота системы собственных функций}

\begin{mytheorem}
Пусть
\[
f\in\operatorname{Dom}\mathcal L,
\qquad
c_k:=\int_a^b f(x)y_k(x)\,\dd x.
\]
Тогда
\[
\sum_{k=1}^{\infty}c_ky_k(x)
\]
сходится к $f$ в $L_2(a,b)$, то есть
\[
\int_a^b
\left|
 f(x)-\sum_{k=1}^{N}c_ky_k(x)
\right|^2\dd x
\longrightarrow0.
\]
\end{mytheorem}

\begin{myproof}
Положим
\[
f_N(x)
=
f(x)-\sum_{k=1}^{N-1}c_ky_k(x).
\]
Тогда
\[
f_N\in\operatorname{Dom}\mathcal L
\]
и
\[
\int_a^b f_N(x)y_k(x)\,\dd x=0,
\qquad k=1,\ldots,N-1.
\]
Следовательно,
\[
J[f_N]
\ge
\lambda_N\int_a^b f_N(x)^2\,\dd x.
\]
С другой стороны,
\[
J[f_N]
=
\int_a^b(\mathcal Lf_N)f_N\,\dd x
\le
\left(\int_a^b(\mathcal Lf_N)^2\,\dd x\right)^{1/2}
\left(\int_a^b f_N^2\,\dd x\right)^{1/2}.
\]
Поэтому
\[
\lambda_N\|f_N\|_{L_2}^2
\le
\|\mathcal Lf_N\|_{L_2}\|f_N\|_{L_2},
\]
то есть
\[
\|f_N\|_{L_2}
\le
\frac{\|\mathcal Lf_N\|_{L_2}}{\lambda_N}.
\]

Имеем
\[
\mathcal Lf_N
=
\mathcal Lf-
\sum_{k=1}^{N-1}c_k\lambda_ky_k.
\]
Для $k=1,\ldots,N-1$
\[
(\mathcal Lf_N,y_k)
=(f_N,\mathcal Ly_k)
=(f_N,\lambda_ky_k)
=0.
\]
Поэтому
\[
\|\mathcal Lf\|_{L_2}^2
=
\|\mathcal Lf_N\|_{L_2}^2
+
\left\|
\sum_{k=1}^{N-1}c_k\lambda_ky_k
\right\|_{L_2}^2,
\]
и
\[
\|\mathcal Lf_N\|_{L_2}
\le
\|\mathcal Lf\|_{L_2}.
\]
Следовательно,
\[
\boxed{
\|f_N\|_{L_2(a,b)}
\le
\frac{\|\mathcal Lf\|_{L_2(a,b)}}{\lambda_N}
\longrightarrow0.}
\]
\end{myproof}

\begin{myprop}
Для $f\in\operatorname{Dom}\mathcal L$
\[
\boxed{
\int_a^b f(x)^2\,\dd x
=
\sum_{k=1}^{\infty}c_k^2.}
\]
\end{myprop}

\begin{myproof}
Так как
\[
\int_a^b
\left(
 f(x)-\sum_{k=1}^{N}c_ky_k(x)
\right)^2\dd x
\to0,
\]
то
\[
\int_a^b f(x)^2\,\dd x
-2\sum_{k=1}^{N}c_k\int_a^b f(x)y_k(x)\,\dd x
+
\int_a^b\left(\sum_{k=1}^{N}c_ky_k(x)\right)^2\dd x
\to0.
\]
Следовательно,
\[
\int_a^b f(x)^2\,\dd x
-\sum_{k=1}^{N}c_k^2
\to0.
\]
\end{myproof}

\begin{myremark}
\begin{enumerate}
\item Для любого $f\in L_2(a,b)$
\[
\int_a^b f(x)^2\,\dd x
=
\sum_{k=1}^{\infty}c_k^2,
\]
так как $\operatorname{Dom}\mathcal L$ плотна в $L_2(a,b)$.
\item Если $f\in\operatorname{Dom}\mathcal L$, то
\[
\sum_{k=1}^{\infty}c_ky_k(x)
\]
сходится равномерно к $f(x)$,
\[
y_k(x)=O(1),
\qquad
c_k=O\!\left(\frac1{k^2}\right),
\qquad k\to\infty.
\]
\end{enumerate}
\end{myremark}

\begin{mylemma}
Других собственных чисел у задачи Штурма--Лиувилля нет.
\end{mylemma}

\begin{myproof}
Пусть $y$ --- собственная функция,
\[
\mathcal Ly=\lambda y,
\qquad
\lambda\ne\lambda_k
\quad\forall k.
\]
Так как $\mathcal L$ симметричен,
\[
y\perp y_k
\qquad\forall k,
\]
следовательно,
\[
c_k=0
\qquad\forall k.
\]
По равенству Парсеваля
\[
\int_a^b y(x)^2\,\dd x
=
\sum_{k=1}^{\infty}c_k^2=0,
\]
откуда
\[
y\equiv0,
\]
что невозможно для собственной функции.
\end{myproof}

\begin{myremark}
Итоговая форма для весовой задачи.
Пусть
\[
p\in C^1[a,b],
\qquad
q,r\in C[a,b],
\qquad
p(x)>0,
\quad r(x)>0.
\]
Рассмотрим
\[
J[y]=\int_a^b(py'^2+qy^2)\,\dd x
\]
и множества
\[
M_n=
\left\{
 y\in C^1[a,b]:
 y(a)=y(b)=0,
 \int_a^b y(x)^2r(x)\,\dd x=1,
\right.
\]
\[
\left.
\hspace{35mm}
\int_a^b y(x)y_k(x)r(x)\,\dd x=0,
\quad k=1,\ldots,n-1
\right\}.
\]
Тогда
\[
\lambda_n:=\inf_{M_n}J
\]
достигается на некоторой функции $y_n$, причём
\[
\boxed{
(p y_n')'+(\lambda_nr-q)y_n=0,}
\qquad
\boxed{y_n(a)=y_n(b)=0.}
\]
Других собственных значений и собственных функций нет,
\[
\lambda_n\to+\infty.
\]
Если
\[
c_k:=\int_a^b f(x)y_k(x)r(x)\,\dd x,
\]
то
\[
\sum_{k=1}^{\infty}c_ky_k(x)=f(x)
\]
в соответствующем смысле, и
\[
\boxed{
\int_a^b f(x)^2r(x)\,\dd x
=
\sum_{k=1}^{\infty}c_k^2.}
\]
\end{myremark}

% ============================================================
% §4. Многомерные вариационные задачи
% ============================================================

\section{Многомерные вариационные задачи}

\subsection{Уравнение Эйлера--Остроградского}

Пусть
\[
\Omega\subset\mathbb R^k
\]
--- ограниченная область и
\[
u:\overline\Omega\to\mathbb R^n.
\]
Положим
\[
X=C^1(\overline\Omega,\mathbb R^n),
\]
\[
\|u\|_X
=
\sum_{j=1}^{n}
\left(
\max_{\overline\Omega}|u_j(x)|
+
\max_{\overline\Omega}|\nabla u_j(x)|
\right).
\]
Рассмотрим функционал
\[
J[u]
=
\int_\Omega L(x,u(x),u'(x))\,\dd x,
\]
где
\[
L:\overline\Omega\times\mathbb R^n\times\mathbb R^{kn}\to\mathbb R,
\qquad
(x,u,w)\longmapsto L(x,u,w),
\]
а
\[
u'(x)
=
\left(
\frac{\partial u_j(x)}{\partial x_i}
\right)_{\substack{i=1,\ldots,k\\j=1,\ldots,n}},
\qquad
w_{ij}=\frac{\partial u_j}{\partial x_i}.
\]
Если
\[
L\in C(\overline\Omega\times\mathbb R^n\times\mathbb R^{kn}),
\]
то $J$ непрерывен на $X$.
Если
\[
L\in C^1(\overline\Omega\times\mathbb R^n\times\mathbb R^{kn}),
\]
то
\[
J\in C^1(X)
\]
и
\[
\boxed{
\delta J(u;h)
=
\int_\Omega
\left(
\sum_{j=1}^{n}
\frac{\partial L}{\partial u_j}(x,u,u')h_j(x)
+
\sum_{i=1}^{k}\sum_{j=1}^{n}
\frac{\partial L}{\partial w_{ij}}(x,u,u')
\frac{\partial h_j(x)}{\partial x_i}
\right)\dd x.}
\]

\begin{mytheorem}
Пусть
\[
L\in C^1(\overline\Omega\times\mathbb R^n\times\mathbb R^{kn}),
\]
граница $\partial\Omega$ достаточно гладкая и задана функция
\[
\varphi\in C(\partial\Omega).
\]
Пусть $u_0$ --- локальный экстремум функционала
\[
J[u]=\int_\Omega L(x,u,u')\,\dd x
\]
на множестве
\[
\left\{
 u\in C^1(\overline\Omega,\mathbb R^n):
 u\big|_{\partial\Omega}=\varphi
\right\},
\]
причём
\[
u_0\in C^2(\overline\Omega,\mathbb R^n).
\]
Тогда
\[
\boxed{
\sum_{i=1}^{k}
\frac{\partial}{\partial x_i}
\left(
\frac{\partial L(x,u_0(x),u_0'(x))}{\partial w_{ij}}
\right)
=
\frac{\partial L(x,u_0(x),u_0'(x))}{\partial u_j},
\quad j=1,\ldots,n.}
\]
\end{mytheorem}

\begin{myremark}
\begin{enumerate}
\item Производная $\dfrac{\partial}{\partial x_i}$ в уравнении берётся как полная производная, но только по переменной $x_i$.
\item Система
\[
\begin{cases}
\displaystyle
\sum_{i=1}^{k}
\frac{\partial}{\partial x_i}
\left(
\frac{\partial L}{\partial w_{ij}}
\right)
=
\dfrac{\partial L}{\partial u_j},
& j=1,\ldots,n,\\[3mm]
u_0\big|_{\partial\Omega}=\varphi
\end{cases}
\]
называется системой уравнений Эйлера--Остроградского.
\end{enumerate}
\end{myremark}

\begin{myproof}
Так как $u_0$ --- локальный экстремум, то
\[
\delta J(u_0;h)=0
\]
для всех $h\in X$ таких, что
\[
h\big|_{\partial\Omega}=0.
\]
Следовательно,
\[
\int_\Omega
\left(
\sum_{j=1}^{n}
\frac{\partial L(x,u_0,u_0')}{\partial u_j}h_j(x)
+
\sum_{i=1}^{k}\sum_{j=1}^{n}
\frac{\partial L(x,u_0,u_0')}{\partial w_{ij}}
\frac{\partial h_j(x)}{\partial x_i}
\right)\dd x
=0.
\]
Интегрируя вторую сумму по частям и используя $h|_{\partial\Omega}=0$, получаем
\[
\int_\Omega
\sum_{j=1}^{n}
\left(
\frac{\partial L(x,u_0,u_0')}{\partial u_j}
-
\sum_{i=1}^{k}
\frac{\partial}{\partial x_i}
\left(
\frac{\partial L(x,u_0,u_0')}{\partial w_{ij}}
\right)
\right)h_j(x)\,\dd x
=0.
\]
Фиксируем $j$ и берём
\[
h_1=\cdots=h_{j-1}=h_{j+1}=\cdots=h_n=0.
\]
По основной лемме вариационного исчисления
\[
\frac{\partial L(x,u_0,u_0')}{\partial u_j}
=
\sum_{i=1}^{k}
\frac{\partial}{\partial x_i}
\left(
\frac{\partial L(x,u_0,u_0')}{\partial w_{ij}}
\right),
\qquad
j=1,\ldots,n.
\]
\end{myproof}






\end{document}